\documentclass[12pt,a4paper,times]{article}
\usepackage{fancyvrb}
\usepackage{color}
\usepackage{comment}
\usepackage{pgf}
\usepackage{relsize}
\usepackage{geometry}
\usepackage{fancyhdr}
\usepackage{lastpage}

\pagestyle{fancy}

\geometry{
  %includeheadfoot,
  margin=2.54cm
}
\fvset{tabsize=2}
\fvset{fontsize=\small}

\lhead{}
\chead{}
\rhead{}
\lfoot{SDP Primer}
\cfoot{Page \thepage\ of \pageref{LastPage}}
\rfoot{Exercise sheet 03}
\renewcommand{\headrulewidth}{0pt}
\renewcommand{\footrulewidth}{0.4pt}

\title{The SDP Primer}
\author{Exercise sheet 03}
\date{Session: 2016-17}

% https://www.cis.upenn.edu/~matuszek/cit591-2013/Assignments/00-warm-up-exercises.html
\begin{document}
\maketitle
\thispagestyle{fancy}

%\fontsize{18}{24}
%\selectfont
\section{Purposes of these exercises}\label{purposes-of-this-assignment}
\begin{itemize}
\item
To give you more practice with Scala, methods in particular
\item
To stress the importance of following specifications exactly
\end{itemize}

\section{General idea of the exercises}\label{general-idea-of-the-assignment}

For each of the numbers 1 through 100, determine some characteristics of
those numbers. Print out the numbers, one per line, along with a list of
its characteristics (on the same line).
Your output should look approximately like this:
\begin{Verbatim}
    1   c, h, t, s, not sm, ht, not pr, d
    2   p, u, not t, not s, sm, ht, pr, d
    3   p, u, t, not s, not smug, ht, not pr, d
    4   c, u, not t, s, not sm, ht, not pr, d
    5   p, u, not t, not s, sm, dis, not pr, d
    6   c, u, t, not s, not sm, dis, pr, per
    ...
\end{Verbatim}
where
\begin{Verbatim}
c = composite
d = deficient
dis = dishonest
h = happy
ht = honest
per = perfect
p = prime
pr = pronic
sm = smug
s = square
t = triangular
u = unhappy
\end{Verbatim}

\section{Details}\label{details}

\subsection{General}\label{general}

Create an object \verb!NumberPersonalities! in a file
\verb!NumberPersonalities.scala!.
Within the \verb!NumberPersonalities! object, declare a variable as \verb!val limit = 100!.
Within the \verb!NumberPersonalities! object, write a
\begin{Verbatim}
def main(args: Array[String]) 
\end{Verbatim}
method to test the
numbers \verb!1! through \verb!limit!, inclusive, and print out
the results in the form shown above. In \verb!main!, do \emph{not}
assume that \verb!limit! is 100, but write your program to handle
\emph{any} positive \verb!Int!. This makes it easy to change the limit later.

\verb!main! should probably be the \emph{last} method you write,
after you have written \emph{and tested} all the methods that it will
need to call.

A \emph{predicate} is a method that returns a \texttt{Boolean} value
(\texttt{true} or \texttt{false}). Within the
\verb!NumberPersonalities! object, write the following predicates:
\begin{itemize}
\item
  \texttt{isPrime(n:\ Int):\ Boolean}
\item
  \texttt{isHappy(n:\ Int):\ Boolean}
\item
  \texttt{isTriangular(n:\ Int):\ Boolean}
\item
  \texttt{isSquare(n:\ Int):\ Boolean}
\item
  \texttt{isSmug(n:\ Int):\ Boolean}
\item
  \texttt{isHonest(n:\ Int):\ Boolean}
\item
  \texttt{isPronic(n:\ Int):\ Boolean}
\item
  \texttt{isDeficient(n:\ Int):\ Boolean}
\item
  \texttt{isPerfect(n:\ Int):\ Boolean}
\item
  \texttt{isAbundant(n:\ Int):\ Boolean}
\end{itemize}
And also write the following method:
\begin{itemize}
\item
  \texttt{sumOfPositiveDivisorsOf(n:\ Int):\ Int}
\end{itemize}
You should provide tests for all your methods.
\begin{itemize}
\item
  Testing does \emph{not} depend on the parameter names (\texttt{n} in
  each of the above). You can use whatever parameter names you like. For
  these methods, however, \texttt{n} is as good a name as any.
\item
  In most cases, Scala can deduce the return type (\texttt{Boolean} in
  all but one of the above methods), and you don't have to specify it.
\end{itemize}

\subsection{Prime and composite numbers}\label{prime-and-composite-numbers}

A positive number is \textbf{prime} if its only positive divisors are
itself and one. For example, \verb!7! is prime because it is evenly
divisible by \verb!1! and \verb!7!, but not by any number in
between (in particular, it is not evenly divisible
by \verb!2!, \verb!3!, \verb!4!, \verb!5!, or \verb!6!).

A positive number is \textbf{composite} if it is not prime. For
example, \verb!10! is composite because it is evenly divisible
by \verb!2! and \verb!5!; \verb!12! is composite because it is evenly divisible
by \verb!2!, \verb!3!, \verb!4!, and \verb!6!. As a
special case, \verb!1! is considered to be composite.

\begin{quote}
\textbf{Programming note.} You can test whether a number \verb!n! is
prime by dividing it by each of the numbers \texttt{2} through
\texttt{n-1}, and testing if the remainder is zero (use the
\textbf{\texttt{\%}} operator). A zero remainder means \verb!n! was
evenly divisible by that divisor, hence is composite.
\end{quote}

\subsection{Happy and unhappy numbers}\label{happy-and-unhappy-numbers}

Repeatedly apply the following procedure to a number:
\begin{enumerate}
\item Square each of the digits of the number.
\item Add the squares together.
\end{enumerate}
If by doing this you eventually get to \verb!1!, then the number is \textbf{happy}.
For example, if you start with \verb!19!:
\begin{itemize}
\item
  \texttt{12\ +\ 92\ =\ 1\ +\ 81\ =\ 82}
\item
  \texttt{82\ +\ 22\ =\ 64\ +\ 4\ =\ 68}
\item
  \texttt{62\ +\ 82\ =\ 36\ +\ 64\ =\ 100}
\item
  \texttt{12\ +\ 02\ +\ 02\ =\ 1\ +\ 0\ +\ 0\ =\ 1} and the number is happy.
\end{itemize}
If instead you get into an infinite loop, the number is
\textbf{unhappy}. So if your program runs forever, the number is
unhappy. This isn't a very useful test, however. Fortunately, every
unhappy number will eventually get into this cycle:
\begin{quote}
\texttt{4,\ 16,\ 37,\ 58,\ 89,\ 145,\ 42,\ 20,\ 4,\ ...}
\end{quote}
so if you get any one of these numbers, say \verb!4!, then you can
stop and conclude that the number is unhappy.
\begin{quote}
\textbf{Programming note.} You can get the last digit of a
number \verb!n! with the expression \texttt{(n\ \%\ 10)}. You can
subsequently discard this digit by doing an ``integer division'' by 10;
when you get a zero, there are no more digits.
\begin{quote}
Example: Taking the number \verb!1024!:
\begin{quote}
\texttt{~~~~~~~~~~~~~~~~1024\ \ \ \ \ \ \ \ \ \ \%\ 10\ =\ 4\ \ \ \ \ \ \ \ \ 1024\ /\ 10\ =\ 102,\ 102\ \%\ 10\ =\ 2\ \ \ \ \ \ \ \ \ ~102\ /\ 10\ =\ 10,~~\ 10\ \%\ 10\ =\ 0\ \ \ \ \ \ \ \ \ ~~10\ /\ 10\ =\ 1,~~~~\ 1\ \%\ 10\ =\ 1\ \ \ \ \ \ \ \ \ ~~\ 1\ /\ 10\ =\ 0\ }(stop)
\end{quote}
\end{quote}
\end{quote}

\subsection{Triangular numbers}\label{triangular-numbers}

\centerline{\includegraphics{images/triangular_numbers}}

A \textbf{triangular number} is a number of the form \texttt{1~+~2~+~3~+~...~+~n}, for some
positive\texttt{\ n}. These numbers are called triangular because, if
you have that many objects, you can arrange them in an equilateral
triangle (see figure).
The first few triangular numbers are
\texttt{1,\ 3,\ 6,\ 10,\ 15}\ldots{}.
\begin{quote}
\textbf{Programming note.} You can test if a number is triangular by
subtracting the number 1 from it, then subtracting 2, then subtracting
3, etc. If at some point you get zero, then the starting number is
triangular; but if you first get a negative number, then the number is
not triangular.

Example:

\noindent\texttt{6\ -\ 1\ =\ 5\ ~~~~5\ -\ 2\ =\ 3\ ~~~~3\ -\ 3\ =\ 0\ }and
the number is triangular.
\end{quote}

\subsection{Square numbers}\label{square-numbers}

A \textbf{square number}

\centerline{\includegraphics[width=0.86458in,height=1.06250in]{images/square-numbers}}

\noindent is a number of the form \texttt{1~+~3~+~5~+~7~+~...~+~n}, for some odd
positive\texttt{\ n}. (The figure should help you understand why this
definition is the same as the usual definition of square numbers.)

\begin{quote}
\textbf{Programming note.} To test if a number is \textbf{square}, do
the following. For each of the odd numbers \texttt{1~+~3~+~5~+~7~+~...},
subtract that number, and see if you get zero (square) or a negative
number (not square).
\end{quote}

\subsection{Smug numbers}\label{smug-numbers}

A number is \textbf{smug} if it is the sum of two square numbers. The
first few smug numbers are \texttt{2} (1+1), \texttt{5} (1+4),
\texttt{8} (4+4), \texttt{10} (1+9), \texttt{13} (4+9), etc.

\begin{quote}
\textbf{Programming note.} To test if a number is smug, do the
following. For each of the square numbers (which you can get as
\texttt{1*1,\ 2*2,\ 3*3}, etc.), subtract that number, and see if the
remainder is also a square number. You can stop if the remainder is less
than the square number you subtracted.
\end{quote}

\subsection{Honest and dishonest
numbers}\label{honest-and-dishonest-numbers}

A number is \textbf{dishonest} if it ``pretends'' to be a square number,
but isn't one. Specifically, n is dishonest if there is a
number\texttt{\ k\ }such that\texttt{\ n/k~=~k\ }(using integer
division), but\texttt{\ \ \ \ k*k\ }is not\texttt{\ n}. A number is
\textbf{honest} if it is not dishonest.

\begin{quote}
\textbf{Programming note.} To find\texttt{\ k}, try the numbers 1, 2, 3,
\ldots{} until you find a\texttt{\ k\ }such that either \texttt{n/k~=~k}
or \texttt{n/k~\textless{}~k}.
\end{quote}

\subsection{Pronic numbers}\label{pronic-numbers}

A number is \textbf{pronic} if it is the product of two consecutive
numbers. The first few pronic numbers are \texttt{2} (1*2), \texttt{6}
(2*3), \texttt{12} (3*4), \texttt{20} (4*5), etc.

\begin{quote}
\textbf{Programming note.} Starting with\texttt{\ k\ =\ 1}, compute
\texttt{k\ *\ (k\ +\ 1)} until the result either equals or exceeds the
number you are testing.
\end{quote}

\subsection{Deficient, perfect, and abundant
numbers}\label{deficient-perfect-and-abundant-numbers}

Every positive number is either \textbf{deficient}, \textbf{perfect}, or
\textbf{abundant}. To determine which it is, add up all the positive
numbers (including \texttt{1}) that divide it evenly.

\begin{itemize}
\item
  If the sum of the divisors is less than the number, the number is
  \textbf{deficient}.
\item
  If the sum of the divisors exactly equals the number, the number is
  \textbf{perfect}.
\item
  If the sum of the divisors is greater than the number, the number is
  \textbf{abundant}.
\end{itemize}
For example, \texttt{24} is evenly divisible by \texttt{1}, \texttt{2},
\texttt{3}, \texttt{4}, \texttt{6}, \texttt{8}, and \texttt{12}. Because
\texttt{1+2+3+4+6+8+12\ =\ 36}, and \texttt{36\ \textgreater{}\ 24}, the
number \texttt{24} is abundant.

A fundamental rule of good programming is the \textbf{DRY Principle}:
\emph{Don't Repeat Yourself!} In each of three methods
(\texttt{isDeficient}, \texttt{isPerfect}, and \texttt{isAbundant}), you
need to compute the sum of the divisors of a number. It's poor style to
write this same code in all three methods; instead, write a
\texttt{sumOfPositiveDivisorsOf} method that each of the three methods
can use.
\end{document}